% GENERATED FILE. Edit canonical lessons, then run npm run generate:books.
% canonical-source-hash: fnv1a64:25cc6154fe438305
% canonical-lessons: MW-C44-pag, MW-C44-matho, MW-C44-bal, MW-C44-dant, MW-C44-galo

\chapter{Foot, Head, Hair, Tooth, Throat}
\label{ch:mw-foot-head-hair-tooth-throat}
\hlchaptermodality{44}

\begin{chapteropening}
\textbf{By the end of this chapter:} \emph{I can say a foot, the head, hair, a tooth and the throat.}

Complete the last lesson of chapter 44: I can say a foot, the head, hair, a tooth and the throat.
\end{chapteropening}

\section[\texorpdfstring{\mw{पग} (pag)}{पग (pag)}]{\mw{पग} (pag) — a foot}

\label{lesson:MW-C44-pag}



\begin{quote}
Before the new one: say the Marwadi for nineteen, then the Marwadi for no.
\end{quote}

\subsection*{You'll want to know: \mw{पग}}
\textbf{\mw{पग}} — \emph{pag} — \textquotedblleft{}a foot\textquotedblright{}.

\begin{cousinweb}
\textbf{\mw{पग}} (\emph{pag}) is a tadbhava, a word inherited from Sanskrit and reshaped by sound change. It comes from \emph{padga}, \textquotedblleft{}afoot; footsoldier\textquotedblright{}. CDIAL lists Old Marwari \emph{paga} under it (CDIAL 7766). It is built on \emph{pad}, \textquotedblleft{}foot\textquotedblright{}, the same ancient word as English \emph{foot}.
\end{cousinweb}

\begin{grammarlens}[title={What you've built}]
One more word for this chapter.
\end{grammarlens}

\subsection*{Your turn}
\begin{itemize}
\raggedright
  \item \emph{Say it:} \emph{pag}
  \item \emph{Say it:} \emph{pag}, once more
  \item \emph{Say it:} say \emph{pag}
  \item \emph{Recall:} say the Marwadi for nineteen, then the Marwadi for no, then say \emph{pag} again
\end{itemize}

\begin{tcolorbox}[breakable,colback=teal!4,colframe=teal!35!black,title={Before you move on}]
What does \mw{पग} mean? (\textquotedblleft{}A foot\textquotedblright{}.) Say it once more.
\end{tcolorbox}
\section[\texorpdfstring{\mw{माथो} (mātho)}{माथो (mātho)}]{\mw{माथो} (mātho) — the head}

\label{lesson:MW-C44-matho}



\begin{quote}
Before the new one: say the Marwadi for no, then the Marwadi for a foot.
\end{quote}

\subsection*{You'll want to know: \mw{माथो}}
\textbf{\mw{माथो}} — \emph{mātho} — \textquotedblleft{}the head\textquotedblright{}.

\begin{cousinweb}
\textbf{\mw{माथो}} (\emph{mātho}) is a tadbhava, a word inherited from Sanskrit and reshaped by sound change. It comes from \emph{masta}, \textquotedblleft{}head\textquotedblright{}. CDIAL lists Marwari \emph{mātho} under it (CDIAL 9926). Sanskrit also had \emph{mastaka}, \textquotedblleft{}head, skull\textquotedblright{}.
\end{cousinweb}

\begin{grammarlens}[title={What you've built}]
One more word for this chapter.
\end{grammarlens}

\subsection*{Your turn}
\begin{itemize}
\raggedright
  \item \emph{Say it:} \emph{mātho}
  \item \emph{Say it:} \emph{mātho}, once more
  \item \emph{Say it:} say \emph{mātho}
  \item \emph{Recall:} say the Marwadi for no, then the Marwadi for a foot, then say \emph{mātho} again
\end{itemize}

\begin{tcolorbox}[breakable,colback=teal!4,colframe=teal!35!black,title={Before you move on}]
What does \mw{माथो} mean? (\textquotedblleft{}The head\textquotedblright{}.) Say it once more.
\end{tcolorbox}
\section[\texorpdfstring{\mw{बाळ} (bāḷ)}{बाळ (bāḷ)}]{\mw{बाळ} (bāḷ) — hair}

\label{lesson:MW-C44-bal}



\begin{quote}
Before the new one: say the Marwadi for a foot, then the Marwadi for the head.
\end{quote}

\subsection*{You'll want to know: \mw{बाळ}}
\textbf{\mw{बाळ}} — \emph{bāḷ} — \textquotedblleft{}hair\textquotedblright{}.

\begin{cousinweb}
\textbf{\mw{बाळ}} (\emph{bāḷ}) is a tadbhava, a word inherited from Sanskrit and reshaped by sound change. It comes from \emph{vāla}, \textquotedblleft{}hair of tail, tail, hair\textquotedblright{}. CDIAL gives no Marwari form in this entry (CDIAL 11572); the match is through Hindi \emph{bāl} and Gujarati \emph{vāḷ}. Marwadi, like Gujarati, says the \emph{l} curled back.
\end{cousinweb}

\begin{grammarlens}[title={What you've built}]
One more word for this chapter.
\end{grammarlens}

\subsection*{Your turn}
\begin{itemize}
\raggedright
  \item \emph{Say it:} \emph{bāḷ}
  \item \emph{Say it:} \emph{bāḷ}, once more
  \item \emph{Say it:} say \emph{bāḷ}
  \item \emph{Recall:} say the Marwadi for a foot, then the Marwadi for the head, then say \emph{bāḷ} again
\end{itemize}

\begin{tcolorbox}[breakable,colback=teal!4,colframe=teal!35!black,title={Before you move on}]
What does \mw{बाळ} mean? (\textquotedblleft{}Hair\textquotedblright{}.) Say it once more.
\end{tcolorbox}
\section[\texorpdfstring{\mw{दांत} (dānt)}{दांत (dānt)}]{\mw{दांत} (dānt) — a tooth}

\label{lesson:MW-C44-dant}



\begin{quote}
Before the new one: say the Marwadi for the head, then the Marwadi for hair.
\end{quote}

\subsection*{You'll want to know: \mw{दांत}}
\textbf{\mw{दांत}} — \emph{dānt} — \textquotedblleft{}a tooth\textquotedblright{}.

\begin{cousinweb}
\textbf{\mw{दांत}} (\emph{dānt}) is a tadbhava, a word inherited from Sanskrit and reshaped by sound change. It comes from \emph{danta}, \textquotedblleft{}tooth\textquotedblright{}. CDIAL lists Marwari \emph{dā̃t} under it (CDIAL 6152). English \emph{tooth} and Latin \emph{dens} are the same ancient word.
\end{cousinweb}

\begin{grammarlens}[title={What you've built}]
One more word for this chapter.
\end{grammarlens}

\subsection*{Your turn}
\begin{itemize}
\raggedright
  \item \emph{Say it:} \emph{dānt}
  \item \emph{Say it:} \emph{dānt}, once more
  \item \emph{Say it:} say \emph{dānt}
  \item \emph{Recall:} say the Marwadi for the head, then the Marwadi for hair, then say \emph{dānt} again
\end{itemize}

\begin{tcolorbox}[breakable,colback=teal!4,colframe=teal!35!black,title={Before you move on}]
What does \mw{दांत} mean? (\textquotedblleft{}A tooth\textquotedblright{}.) Say it once more.
\end{tcolorbox}
\section[\texorpdfstring{\mw{गळो} (gaḷo)}{गळो (gaḷo)}]{\mw{गळो} (gaḷo) — the throat}

\label{lesson:MW-C44-galo}



\begin{quote}
Before the new one: say the Marwadi for hair, then the Marwadi for a tooth.
\end{quote}

\subsection*{You'll want to know: \mw{गळो}}
\textbf{\mw{गळो}} — \emph{gaḷo} — \textquotedblleft{}the throat\textquotedblright{}.

\begin{cousinweb}
\textbf{\mw{गळो}} (\emph{gaḷo}) is a tadbhava, a word inherited from Sanskrit and reshaped by sound change. It comes from \emph{gala}, \textquotedblleft{}throat, neck\textquotedblright{}. CDIAL lists Marwari \emph{gaḷo} under it (CDIAL 4070).
\end{cousinweb}

\begin{grammarlens}[title={What you've built}]
One more word for this chapter.
\end{grammarlens}

\subsection*{Your turn}
\begin{itemize}
\raggedright
  \item \emph{Say it:} \emph{gaḷo}
  \item \emph{Say it:} \emph{gaḷo}, once more
  \item \emph{Say it:} say \emph{gaḷo}
  \item \emph{Recall:} say the Marwadi for hair, then the Marwadi for a tooth, then say \emph{gaḷo} again
\end{itemize}

\begin{tcolorbox}[breakable,colback=teal!4,colframe=teal!35!black,title={Before you move on}]
What does \mw{गळो} mean? (\textquotedblleft{}The throat\textquotedblright{}.) Say it once more.
\end{tcolorbox}
